\documentclass{article}
\usepackage{PRIMEarxiv} % Core package for the PrimeAI Template
\usepackage{amsmath, amssymb, amsthm, bm, dcolumn} % Add amsthm here for the proof environment
\usepackage[numbers,sort&compress]{natbib} % Natbib for citations
\usepackage{graphicx} % For high-quality images
\usepackage{multicol} % Multi-column figures/tables
\usepackage{hyperref} % Hyperlinks
\usepackage{listings} % Code listings
\usepackage{authblk} % For structured affiliations
% Define theorem style (optional)
\theoremstyle{plain}
\newtheorem{theorem}{Theorem}
\renewcommand{\qedsymbol}{}

% Custom commands for primes and qubit encoding
\newcommand{\primeQubit}[1]{|p_{#1}\rangle}
\newcommand{\primeGate}[1]{U_{p_{#1}}}

\usepackage{wrapfig}
\usepackage[pscoord]{eso-pic}
\usepackage[fulladjust]{marginnote}
\reversemarginpar

% Typesetting improvements without footnote patching
\usepackage[protrusion=true, expansion=true, tracking=false]{microtype}
\microtypecontext{spacing=nonfrench}

% Line numbers
\usepackage[right]{lineno}

% Text layout - adjust as needed
\raggedright
\setlength{\parindent}{0.5cm}
\textwidth 5.25in 
\textheight 8.75in

% Set double spacing
\usepackage{setspace} 
\doublespacing

% Adjust width for specific content
\usepackage{changepage}

% Adjust caption style
\usepackage[aboveskip=1pt,labelfont=bf,labelsep=period,singlelinecheck=off]{caption}

% Remove brackets from references
\makeatletter
\renewcommand{\@biblabel}[1]{\quad#1.}
\makeatother

% Header, footer, and page numbers
\usepackage{lastpage,fancyhdr}
\pagestyle{fancy}
\fancyhf{}
\fancyhead[L]{Preprint - PrimeAI Enhanced Template}
\fancyfoot[C]{\scriptsize Multiplicity Theory © 2024 Ryan Van Gelder - Citizen Gardens \\ Licensed Under MIT and CC BY-NC-SA 4.0.}
\fancyfoot[R]{Page \thepage\ of \pageref{LastPage}}
\renewcommand{\footrule}{\hrule height 2pt \vspace{2mm}}

\begin{document}

\title{Multiplicity \& Stirling-Ramanujan Constants}

\author{Ryan O. Van Gelder}
\affil{Citizen Gardens - The Foundation of Multiplicity \\ \texttt{info@citizengardens.org}}
\date{\today}

\maketitle

\begin{abstract}
This article explores the integration of Multiplicity Theory and Stirling-Ramanujan constants. Leveraging concepts from prime-based encoding, tensor networks, and resummation techniques, Multiplicity efficiently handles divergent series, quantum states, and large datasets. The connections to foundational works, such as Einstein's General Relativity \cite{Einstein1916} and Ramanujan's resummation techniques \cite{Ramanujan1957}, highlight the theoretical underpinnings and practical applications across quantum mechanics, astrophysics, and mathematical physics.
\end{abstract}

\section{Introduction}
Multiplicity Theory provides a powerful framework for analyzing complex systems, bridging general relativity, quantum mechanics, and string theory \cite{Witten1995, Green1987}. This article focuses on how the Stirling-Ramanujan constants integrate into MCP’s mathematical structures, enabling precise simulations and computations in quantum systems.

Stirling-Ramanujan constants, arising in asymptotic expansions and divergent series, play a crucial role in resummation techniques. Foundational contributions from Hardy \cite{Hardy1949} and Ramanujan \cite{Ramanujan1957} emphasize the importance of these constants in mathematical and physical systems.

\section{Mathematical Framework of Multiplicity}
Multiplicity Theory employs eigenvalue-based modeling to unify quantum and classical systems. For example, a time-dependent multiplicity equation:
\begin{equation}
M(t) = \sum_{i=1}^N \lambda_i \mu_i \cdot e^{i \theta_i(t)} \cdot v_i,
\end{equation}
extends classical eigenvalue theory into quantum domains \cite{Schrodinger1926, Maldacena1999}. Stirling-Ramanujan constants integrate naturally into this framework, offering compact representations of divergent series.

The Stirling approximation is given by:
\begin{equation}
n! \sim \sqrt{2 \pi n} \left(\frac{n}{e}\right)^n e^{\frac{B_1}{12n}},
\end{equation}
where $B_1$ is the first Bernoulli number. This approximation is crucial for modeling quantum entropy and black hole thermodynamics \cite{Bekenstein1973, Hawking1975}.

\section{Prime-Based Encoding in MCP}
MCP employs prime-based encoding to represent quantum states and mathematical constructs efficiently. By mapping the Stirling-Ramanujan constants to prime factorization schemes, MCP handles large sums and divergent series with ease:
\begin{equation}
S_n = (-1)^{n+1} n! \int_0^{\infty} \frac{1}{p_n(t)} \left( \sum_{k=-1}^n \frac{1}{p_k(t)} - r_n p_{n+1}(t) \right) e^{-t} \frac{dt}{t}.
\end{equation}
Here, $p_n(t)$ are prime-encoded functions representing Stirling-Ramanujan integrals. This encoding enhances MCP's ability to simulate multi-dimensional quantum systems.

\section{Tensor Networks in MCP}
Tensor networks represent quantum states and interactions in MCP, integrating Stirling-Ramanujan constants for modeling wavefunctions and asymptotic expansions:
\begin{equation}
\Phi(t) = \sum_{i,j,k} T_{i,j,k} \Psi_i(t) \otimes \Psi_j(t) \otimes S_k,
\end{equation}
where $T_{i,j,k}$ are tensor coefficients, and $S_k$ are Stirling-Ramanujan constants. This representation allows MCP to compute quantum coherence and entanglement efficiently \cite{Heisenberg1927, Feynman1965}.

\section{Resummation Techniques in Quantum Algorithms}
Resummation techniques are critical in quantum algorithms for handling divergent series. Using the Euler-McLaurin formula, MCP employs Stirling-Ramanujan constants to optimize quantum computations:
\begin{equation}
\Psi(t) = \sum_{n=0}^\infty \alpha_n S_n \cdot e^{i \theta_n(t)},
\end{equation}
where $\alpha_n$ are quantum amplitudes, $S_n$ encode resummation, and $\theta_n(t)$ denotes phase evolution. This formulation stabilizes computations in quantum field theory and string theory \cite{Kaluza1921, Rovelli2004}.

\section{Applications in MCP}

\subsection{Quantum Field Theory Simulations}
Stirling-Ramanujan constants regularize divergences in perturbation expansions, aiding renormalization. For example, the Glaisher-Kinkelin constant, $S_1$, appears in determinant calculations of the Laplace operator, a key component in field theory \cite{Hardy1949}.

\subsection{Thermodynamics and Black Hole Physics}
In thermodynamics, MCP models phase transitions and entropy:
\begin{equation}
S = \frac{k_B A}{4 G} + \beta \chi(M),
\end{equation}
where $\chi(M)$ is the Euler characteristic. These constants stabilize large-scale simulations in astrophysics and cosmology \cite{Bekenstein1973, Hawking1975}.

\subsection{Number Theory and Zeta Functions}
MCP employs Stirling-Ramanujan constants in computing zeta-regularized determinants. The prime-based encoding aligns naturally with the distribution of prime numbers and zeta function expansions \cite{Lagarias2013}.

\section{Conclusion}
The integration of Stirling-Ramanujan constants into MCP demonstrates a powerful synergy between mathematical physics and quantum computation. By leveraging prime-based encoding, tensor networks, and resummation techniques, MCP provides a robust framework for handling complex systems across multiple disciplines.

\bibliographystyle{plain}
\bibliography{references}

\end{document}
